% GENERATED FILE. Edit canonical lessons, then run npm run generate:books.
% canonical-source-hash: fnv1a64:39e725917d9cf898
% canonical-lessons: FR-C497-canne-a-peche, FR-C497-hamecon, FR-C497-appat, FR-C497-pecher, FR-C497-pecheur, FR-R497-revision-490-497

\chapter{Une journée de pêche}
\label{ch:fr-une-journee-de-peche}
\hlchaptermodality{497}

\begin{chapteropening}
\textbf{By the end of this chapter:} \emph{I can talk about a day out fishing.}

Complete the last lesson of chapter 497: I can talk about a day out fishing.
\end{chapteropening}

\section[la canne à pêche]{la canne à pêche — fishing rod}

\label{lesson:FR-C497-canne-a-peche}



\begin{quote}
Before the new one: say the French for to embroider, then the French for sewing.
\end{quote}

\subsection*{You'll want to know: la canne à pêche}
\textbf{la canne à pêche} — \textquotedblleft{}fishing rod\textquotedblright{}.

\begin{quote}
Mon oncle m'a offert une canne à pêche.
\end{quote}

\begin{quote}
\emph{My uncle gave me a fishing rod.}
\end{quote}

\emph{La canne} alone is a walking stick.

\begin{cousinweb}
\emph{la canne à pêche} joins \emph{canne} \textquotedblleft{}cane, rod\textquotedblright{}, from Latin \emph{canna} \textquotedblleft{}a reed\textquotedblright{}, and \emph{pêche} \textquotedblleft{}fishing\textquotedblright{}, built on \emph{pêcher}, from Latin \emph{piscārī} \textquotedblleft{}to fish\textquotedblright{} (Lewis \& Short).
\end{cousinweb}

\begin{grammarlens}[title={What you've built}]
One more word for telling, explaining and describing at B1.
\end{grammarlens}

\subsection*{Your turn}
\begin{itemize}
\raggedright
  \item \emph{Say it:} \emph{la canne à pêche}
  \item \emph{Say it:} \emph{la canne à pêche}, once more
  \item \emph{Say it:} a sentence of your own with \emph{la canne à pêche}
  \item \emph{Recall:} say the French for to embroider, then the French for sewing, then say \emph{la canne à pêche} again
\end{itemize}

\begin{tcolorbox}[breakable,colback=teal!4,colframe=teal!35!black,title={Before you move on}]
What does la canne à pêche mean? (\textquotedblleft{}Fishing rod\textquotedblright{}.) Say it once more.
\end{tcolorbox}
\section[l'hameçon]{l'hameçon — fish hook}

\label{lesson:FR-C497-hamecon}



\begin{quote}
Before the new one: say the French for sewing, then the French for fishing rod.
\end{quote}

\subsection*{You'll want to know: l'hameçon}
\textbf{l'hameçon} — \textquotedblleft{}fish hook\textquotedblright{}.

\begin{quote}
Le poisson a mordu à l'hameçon.
\end{quote}

\begin{quote}
\emph{The fish took the hook.}
\end{quote}

Masculine. \emph{Mordre à l'hameçon} is also to fall for a trick.

\begin{cousinweb}
\emph{l'hameçon} comes from Old French \emph{ameçon}, built on a Vulgar Latin form of Latin \emph{hāmus} \textquotedblleft{}a hook\textquotedblright{} (Lewis \& Short).
\end{cousinweb}

\begin{grammarlens}[title={What you've built}]
One more word for telling, explaining and describing at B1.
\end{grammarlens}

\subsection*{Your turn}
\begin{itemize}
\raggedright
  \item \emph{Say it:} \emph{l'hameçon}
  \item \emph{Say it:} \emph{l'hameçon}, once more
  \item \emph{Say it:} a sentence of your own with \emph{l'hameçon}
  \item \emph{Recall:} say the French for sewing, then the French for fishing rod, then say \emph{l'hameçon} again
\end{itemize}

\begin{tcolorbox}[breakable,colback=teal!4,colframe=teal!35!black,title={Before you move on}]
What does l'hameçon mean? (\textquotedblleft{}Fish hook\textquotedblright{}.) Say it once more.
\end{tcolorbox}
\section[l'appât]{l'appât — bait}

\label{lesson:FR-C497-appat}



\begin{quote}
Before the new one: say the French for fishing rod, then the French for fish hook.
\end{quote}

\subsection*{You'll want to know: l'appât}
\textbf{l'appât} — \textquotedblleft{}bait\textquotedblright{}.

\begin{quote}
Il utilise des vers comme appât.
\end{quote}

\begin{quote}
\emph{He uses worms as bait.}
\end{quote}

Masculine. The verb is \emph{appâter} \textquotedblleft{}to lure\textquotedblright{}.

\begin{cousinweb}
\emph{l'appât} is built on \emph{appâter} (EtymDB), from \emph{à} and Old French \emph{past} \textquotedblleft{}food\textquotedblright{}, which goes back to Latin \emph{pāstus} \textquotedblleft{}feeding\textquotedblright{}, from \emph{pāscere} \textquotedblleft{}to feed\textquotedblright{} (Lewis \& Short).
\end{cousinweb}

\begin{grammarlens}[title={What you've built}]
One more word for telling, explaining and describing at B1.
\end{grammarlens}

\subsection*{Your turn}
\begin{itemize}
\raggedright
  \item \emph{Say it:} \emph{l'appât}
  \item \emph{Say it:} \emph{l'appât}, once more
  \item \emph{Say it:} a sentence of your own with \emph{l'appât}
  \item \emph{Recall:} say the French for fishing rod, then the French for fish hook, then say \emph{l'appât} again
\end{itemize}

\begin{tcolorbox}[breakable,colback=teal!4,colframe=teal!35!black,title={Before you move on}]
What does l'appât mean? (\textquotedblleft{}Bait\textquotedblright{}.) Say it once more.
\end{tcolorbox}
\section[pêcher]{pêcher — to fish; to catch (fish)}

\label{lesson:FR-C497-pecher}



\begin{quote}
Before the new one: say the French for fish hook, then the French for bait.
\end{quote}

\subsection*{You'll want to know: pêcher}
\textbf{pêcher} — \textquotedblleft{}to fish; to catch (fish)\textquotedblright{}.

\begin{quote}
Le dimanche, nous allons pêcher au bord du lac.
\end{quote}

\begin{quote}
\emph{On Sundays we go fishing by the lake.}
\end{quote}

Do not mix it up with \emph{pécher} \textquotedblleft{}to sin\textquotedblright{} or \emph{le pêcher}, the peach tree.

\begin{cousinweb}
\emph{pêcher} is inherited, through Old French \emph{peschier}, from Vulgar Latin \emph{piscāre} (EtymDB), a form of Latin \emph{piscārī} \textquotedblleft{}to fish\textquotedblright{}, from \emph{piscis} \textquotedblleft{}fish\textquotedblright{} (Lewis \& Short).
\end{cousinweb}

\begin{grammarlens}[title={What you've built}]
One more word for telling, explaining and describing at B1.
\end{grammarlens}

\subsection*{Your turn}
\begin{itemize}
\raggedright
  \item \emph{Say it:} \emph{pêcher}
  \item \emph{Say it:} \emph{pêcher}, once more
  \item \emph{Say it:} a sentence of your own with \emph{pêcher}
  \item \emph{Recall:} say the French for fish hook, then the French for bait, then say \emph{pêcher} again
\end{itemize}

\begin{tcolorbox}[breakable,colback=teal!4,colframe=teal!35!black,title={Before you move on}]
What does pêcher mean? (\textquotedblleft{}To fish; to catch (fish)\textquotedblright{}.) Say it once more.
\end{tcolorbox}
\section[le pêcheur]{le pêcheur — fisherman, angler}

\label{lesson:FR-C497-pecheur}



\begin{quote}
Before the new one: say the French for bait, then the French for to fish.
\end{quote}

\subsection*{You'll want to know: le pêcheur}
\textbf{le pêcheur} — \textquotedblleft{}fisherman, angler\textquotedblright{}.

\begin{quote}
Le pêcheur attend depuis le lever du soleil.
\end{quote}

\begin{quote}
\emph{The angler has been waiting since sunrise.}
\end{quote}

A woman is \emph{la pêcheuse}. Do not mix it up with \emph{le pécheur} \textquotedblleft{}sinner\textquotedblright{}.

\begin{cousinweb}
\emph{le pêcheur} is inherited from Latin \emph{piscātor} \textquotedblleft{}a fisherman\textquotedblright{}, from \emph{piscārī} \textquotedblleft{}to fish\textquotedblright{}, from \emph{piscis} \textquotedblleft{}fish\textquotedblright{} (Lewis \& Short).
\end{cousinweb}

\begin{grammarlens}[title={What you've built}]
One more word for telling, explaining and describing at B1.
\end{grammarlens}

\subsection*{Your turn}
\begin{itemize}
\raggedright
  \item \emph{Say it:} \emph{le pêcheur}
  \item \emph{Say it:} \emph{le pêcheur}, once more
  \item \emph{Say it:} a sentence of your own with \emph{le pêcheur}
  \item \emph{Recall:} say the French for bait, then the French for to fish, then say \emph{le pêcheur} again
\end{itemize}

\begin{tcolorbox}[breakable,colback=teal!4,colframe=teal!35!black,title={Before you move on}]
What does le pêcheur mean? (\textquotedblleft{}Fisherman, angler\textquotedblright{}.) Say it once more.
\end{tcolorbox}
\section[(dialogue)]{Retour sur les chapitres 490–497}

\label{lesson:FR-R497-revision-490-497}



\begin{quote}
Say the French for godfather and godmother.
\end{quote}

\subsection*{Your turn}
Say these aloud:

\begin{itemize}
\raggedright
  \item godfather, godmother, godson, eldest, younger
  \item spouse, to separate, separation, to break up, break-up
  \item blended, to adopt, adoption, siblings, large family
  \item pregnancy, to give birth, infant, baby's bottle, to rock
  \item to bring up, to spoil, punishment, to disobey, day nursery
  \item quarrel, to mix up, to deceive, to betray, betrayal
  \item to sew, to knit, ball, to embroider, sewing
  \item fishing rod, fish hook, bait, to fish, fisherman
\end{itemize}

\begin{tcolorbox}[breakable,colback=teal!4,colframe=teal!35!black,title={Before you move on}]
Godfather? (\textbf{le parrain}.) Godmother? (\textbf{la marraine}.) Godson? (\textbf{le filleul}.) Eldest? (\textbf{l'aîné}.) Younger? (\textbf{le cadet}.) Spouse? (\textbf{le conjoint}.) To separate? (\textbf{séparer}.) Separation? (\textbf{la séparation}.) To break up? (\textbf{rompre}.) Break-up? (\textbf{la rupture}.) Blended? (\textbf{recomposé}.) To adopt? (\textbf{adopter}.) Adoption? (\textbf{l'adoption}.) Siblings? (\textbf{la fratrie}.) Large family? (\textbf{la famille nombreuse}.) Pregnancy? (\textbf{la grossesse}.) To give birth? (\textbf{accoucher}.) Infant? (\textbf{le nourrisson}.) Baby's bottle? (\textbf{le biberon}.) To rock? (\textbf{bercer}.) To bring up? (\textbf{élever}.) To spoil? (\textbf{gâter}.) Punishment? (\textbf{la punition}.) To disobey? (\textbf{désobéir}.) Day nursery? (\textbf{la crèche}.) Quarrel? (\textbf{la dispute}.) To mix up? (\textbf{brouiller}.) To deceive? (\textbf{tromper}.) To betray? (\textbf{trahir}.) Betrayal? (\textbf{la trahison}.) To sew? (\textbf{coudre}.) To knit? (\textbf{tricoter}.) Ball? (\textbf{la pelote}.) To embroider? (\textbf{broder}.) Sewing? (\textbf{la couture}.) Fishing rod? (\textbf{la canne à pêche}.) Fish hook? (\textbf{l'hameçon}.) Bait? (\textbf{l'appât}.) To fish? (\textbf{pêcher}.) Fisherman? (\textbf{le pêcheur}.)
\end{tcolorbox}
